图灵机由其指令集决定，指令集是五元组的有限集（即对于每个状态与读入的符号，指定新状态、写出符号以及读写头的移动方式）。五元组的有限集是可数的，因此可以将每个图灵机的指令集与一个数相关联。图灵机的\emph{指标}，就是在固定的此类方案下，与其指令集相关联的数。这样，我们就可以间接地“谈论”图灵机——通过谈论它们的指标。

关于图灵机行为的一个重要问题是：它们最终是否会停机。令 $h(e, n)$ 为这样一个函数：若指标为~$e$ 的图灵机在输入~$n$ 上启动后停机，则函数值为 $1$，否则为 $0$。该函数称为\emph{停机函数}。停机函数本身是否图灵可计算，这一问题称为\emph{停机问题}。答案是否定的：停机问题是不可解的。这是通过对角线论证确立的。

停机问题仅仅是形如“X能否由图灵机完成”的更广泛的一类问题中的一个例子。逻辑的另一个核心问题是\emph{一阶逻辑的判定问题}：是否存在一台图灵机，能够判定给定语句是否有效。这个著名问题同样得到了否定的解答：判定问题是不可解的。这是通过归约论证确立的：对于每台图灵机~$M$ 和输入~$w$，我们可以关联一个一阶语句~$T(M, w) \lif E(M, w)$，该语句有效当且仅当 $M$ 在输入~$w$ 上启动后停机。如果判定问题是可解的，我们就可以用它来解决停机问题。